% ---- pg-494 （书页 482） ----

表现可以用拉亨布鲁克的提留方法来评估. 若 $n_{iM}^{(H)}$ 为第 $i$ 组中错分提留观测的个数, $i=1,2,\cdots,g$, 则期望真实失误率 $E(\mathrm{AER})$ 由下式估计:
\phantomsection\label{eq:11-57}
\begin{equation}
\hat{E}(\mathrm{AER})=\frac{\sum_{i=1}^{g}n_{iM}^{(H)}}{\sum_{i=1}^{g}n_i}
\tag{11-57}
\end{equation}

\noindent\textbf{例 11.12}\quad (用少数几个变量的有效分类)

费希尔(见文献[8])在有关判别函数的开拓性工作中, 对安德森(见文献[1])收集的三种鸢尾属植物的花的数据进行了分析(见练习 11.27, \hyperref[tab:11.5]{表 11.5}). 设三个类别的定义为: $\pi_1$, $\pi_2$, $\pi_3$, 分别代表三种鸢尾属植物. 对每个种的 50 株植物测量了如下四个变量:
\[
\begin{aligned}
X_1&=\text{萼片长度}, & X_2&=\text{萼片宽度}\\
X_3&=\text{花瓣长度}, & X_4&=\text{花瓣宽度}
\end{aligned}
\]
利用\hyperref[tab:11.5]{表 11.5} 中的全部数据, 用线性判别分析得到如下模糊矩阵.

\begin{mvtable}
\begin{tabular}{ll|ccc|c}
\toprule
\multicolumn{6}{c}{预测归属}\\
& & $\pi_1$ & $\pi_2$ & $\pi_3$ & 正确百分率\\
\midrule
\multirow{3}{*}{真实归属} & $\pi_1$ & 50 & 0 & 0 & 100\\
& $\pi_2$ & 0 & 48 & 2 & 96\\
& $\pi_3$ & 0 & 1 & 49 & 98\\
\bottomrule
\end{tabular}
\end{mvtable}

矩阵中的元素用提留方法得出, 故(见\hyperref[eq:11-57]{式(11-57)})
\[
\hat{E}(\mathrm{AER})=3/150=0.02
\]
这个失误率很低, 只有 2\%.

经常可以用较少的变量得到有效的分类. 在实践中可以一次用一个变量, 一次用两个变量, 一次用三个变量等等来进行分类, 然后将这些分类结果与用全部变量分类所得结果进行比较.

如果我们期望 AER 的提留估计值作为准则, 则对鸢尾属植物的数据得出下面的结果.

\begin{mvtable}
\begin{tabular}{l|l}
\toprule
\mvbold{单变量} & \mvbold{错分率}\\
\midrule
$X_1$ & 0.253\\
$X_2$ & 0.480\\
$X_3$ & 0.053\\
$X_4$ & 0.040\\
\midrule
\mvbold{双变量} & \mvbold{错分率}\\
\midrule
$X_1,X_2$ & 0.207\\
$X_1,X_3$ & 0.040\\
$X_1,X_4$ & 0.040\\
$X_2,X_3$ & 0.047\\
$X_2,X_4$ & 0.040\\
$X_3,X_4$ & 0.040\\
\bottomrule
\end{tabular}
\end{mvtable}

我们看到, 单变量 $X_4=$ 花瓣宽度能很好地将三个种分类. 而且, 进一步增加变量几乎得不到任何改进. $X_4$ 的盒形图在\hyperref[fig:11.11]{图 11.11} 中给出. 从图中看得很清楚, 用花瓣宽度可以很好地将三
